\documentclass[11pt]{article}
\usepackage{amsmath, amssymb, amsthm}
\usepackage{geometry}
\geometry{a4paper, margin=1in}
\usepackage{enumitem}
\usepackage{multicol} % Multi-column figures/tables
\usepackage{natbib}
\usepackage{hyperref} % Hyperlinks
\usepackage{listings} % Code listings\
\usepackage[pagewise]{lineno}\linenumbers
\usepackage{authblk} % For structured affiliations
\usepackage{tabularx}
\usepackage{booktabs}
\usepackage{setspace}
\usepackage{lmodern}
\usepackage[T1]{fontenc}
\usepackage{hyperref}
\usepackage{xcolor}
\usepackage{titling}
\usepackage{fancyhdr}
\usepackage{lastpage}
% Define theorem style (optional)
\theoremstyle{plain}
\newtheorem{theorem}{Theorem}
% Theorem environments
\newtheorem{definition}[theorem]{Definition}
\newtheorem{conjecture}[theorem]{Conjecture}
\newtheorem{proposition}[theorem]{Proposition}
\newtheorem{lemma}[theorem]{Lemma}
\nolinenumbers
\hypersetup{
  colorlinks=true,
  linkcolor=black,
  urlcolor=black,
  pdftitle={One-Loop Health-Education System — Build Blueprint},
  pdfauthor={Your Team},
  pdfsubject={Intrinsica, M-Education HEP, Citizen Gardens Overlay},
  pdfkeywords={FHIR, xAPI, ULMS, HIPAA, GDPR, XR, CSL, PER, TCI, MLEM}
}

\begin{document}
\begin{titlepage}
  \centering

  % Optional: left/right logos
  % \begin{minipage}{0.48\textwidth}\flushleft
  %   \includegraphics[height=1.1cm]{logo_left.pdf}
  % \end{minipage}
  % \hfill
  % \begin{minipage}{0.48\textwidth}\flushright
  %   \includegraphics[height=1.1cm]{logo_right.pdf}
  % \end{minipage}

  \vspace*{1cm}

  {\Large\textbf{One-Loop Health-Education System}}\\[6pt]
  {\large Build Blueprint}\\[10pt]
  {\normalsize Intrinsica \(\times\) M-Education HEP \(\times\) Citizen Gardens}\\[20pt]

  \rule{\textwidth}{0.4pt}\\[-4pt]
  \rule{\textwidth}{1.2pt}

  \vspace{1.2cm}

  \begin{tabular}{>{\bfseries}p{3.2cm} p{10.5cm}}
    Document ID: & OLHES-BB-001 \\
    Version:     & v1.1 \\
    Status:      & Build-ready \\
    Date:        & October 15, 2025 \\
    Scope:       & Intrinsica data/AI plane; HEP ULMS/XR; CG overlay \\
    Keywords:    & FHIR R4, HL7 v2, LTI 1.3, xAPI, ULMS, XR, CSL, PER, TCI, MLEM, Ledger, De-ID \\
  \end{tabular}

  \vfill

  \begin{tabular}{>{\bfseries}p{3.2cm} p{10.5cm}}
    Authors:     & \textit{Ryan O. Van Gelder \& Kara Olivarria} \\
    Affiliation: & Citizen Gardens: M-Education Health Nexus \\
    Contact:     & \texttt{info@citizengardens.org} \\
  \end{tabular}

  \vspace{1.2cm}

  \begin{spacing}{1.05}\small
  \textbf{Abstract} \\
  This blueprint specifies an end-to-end, closed-loop architecture linking clinical care and learning. It defines a single data plane (Intrinsica), dual experience skins (patient education | student learning), and shared governance. Interfaces include FHIR/xAPI, ULMS webhooks, and a permissioned ledger with consent-bound provenance.
  \end{spacing}

  \vfill

  \begin{spacing}{1.05}\footnotesize
  \textbf{Notice} \\
  This document contains confidential information belonging to the project partners. Do not distribute, copy, or disclose without prior written permission. All ethics and privacy commitments apply to any derived workproducts.
  \end{spacing}

\end{titlepage}
\newpage
\section{Context and Rationale}
%=============================================================================

\subsection{Problem Statement}
\textbf{Fragmented care.} Clinical, wearable, and behavioral data live in silos; feedback rarely closes the loop back to daily practice. Care plans are static, not signal-responsive. \\
\textbf{Static curricula.} Learning sequences are linear and time-based, not adaptive to readiness, cognitive clarity, or sensory load. \\
\textbf{Low data agency.} Individuals lack granular consent, explainability, and redress. Data reuse is opaque and difficult to audit. \\
\textbf{Inequity.} Neurodiverse learners and underserved communities face pacing mismatches, access barriers, and non-inclusive design.

\subsection{Vision}
\emph{Recursive sovereignty} linking health, learning, and civic capacity: a dual-loop engine where a digital health twin personalizes \emph{what to act on} and a Recursive Assessment Engine personalizes \emph{how capacity is built}, both bounded by a governance tensor that protects dignity, safety, and agency. One role-aware UX integrates wellness practice, mastery growth, and participatory engagement.

\subsection{Objectives}
\begin{enumerate}
  \item \textbf{Outcomes.} Reduce risk trajectories and improve mastery:
    \begin{itemize}
      \item Health: increase intervention adherence; reduce sensor-to-outcome lag $\Delta t_{s\to o}$; flatten adverse risk slope.
      \item Learning: raise $\Delta \mathrm{TCI}$/unit; improve PER uplift and MLEM responsiveness; increase reflection fidelity.
    \end{itemize}
  \item \textbf{Safety.} Enforce pre/post CSL checks with runtime explainability and incident redress; rate-limit intensity by EchoBraid profile.
  \item \textbf{Equity.} Reduce accommodation latency; expand participation diversity; ensure accessible modalities and language support.
  \item \textbf{Sovereignty.} Provide granular, revocable consent; selective disclosure by role; auditable decision traces.
  \item \textbf{Interoperability.} Ingest and emit standards-based resources (e.g., FHIR/HL7); maintain portable schemas for learning and civic events.
\end{enumerate}

\subsection{Guiding Principles}
\begin{description}
  \item[Dignity.] Treat every intervention and learning task as consented, reversible, and paced to human limits.
  \item[Lawfulness.] Bind recommendations to CSL constraints with explainable tensor paths and accountable overrides.
  \item[Reciprocity.] Balance give/receive across community practice; record participation in a transparent ledger.
  \item[Inclusivity.] Design first for neurodiversity and language access; prefer multimodal inputs and outputs.
  \item[Explainability.] Every decision carries sources, bounds, and reasons sufficient for review, audit, and repair.
\end{description}


\newpage


%=============================================================================
\section{Glossary \& Acronyms}
%=============================================================================
\begin{tabularx}{\textwidth}{@{} l X @{}}
\toprule
\textbf{Term} & \textbf{Definition} \\
\midrule
\(\Lambda_m\) & Universal Multiplicity Constant. Alignment parameter that maps health, learning, and civic tensor spaces for cross-domain coherence and traceability. \\
PIRTM & Prime-Indexed Recursive Traceability Model. Identity coherence and lawful provenance across coupled loops. \\
CSL & Conscious Sovereignty Layer. Ethics guardrails for dignity, safety, agency; pre/post checks with explainability and redress. \\
ELM & Reciprocity logic for civic participation used in Citizen Gardens; computes balance of give/receive and quorum dynamics. \\
QARI & Lawful conversational agent for check-ins, reflection coaching, and CSL alignment in dialogue/XR sessions. \\
RAE & Recursive Assessment Engine selecting the next best learning task from signals, mastery, and constraints. \\
PER & Learning-performance metric tracking progress per unit effort and context. \\
MLEM & Micro-Learning Effectiveness Measure estimating effect size of short interventions on mastery and behavior. \\
TCI & Thought Clarity Index approximating cognitive clarity/integration from reflections and task outcomes. \\
EchoBraid & Neurodiversity pacing/modality protocol that rate-limits intensity by sensory profile and context. \\
FHIR/HL7 & Healthcare interoperability standards; FHIR R4 resources used for Patient, Device, Observation, Consent, etc. \\
ULMS & Universal Learning Management System hosting Q-Education, M-Education, and Phenomenal-Edu modules and PD content. \\
\bottomrule
\end{tabularx}


\newpage


\nocite{*}
\bibliographystyle{unsrt}
\bibliography{references}


\end{document}